\section*{Consideraciones generales}

Decisiones y convenciones tomadas a lo largo del documento. Actualizar a medida que surjan nuevas.

\subsection*{Convenciones LaTeX}
\begin{itemize}
    \item Las tablas usan \texttt{[t]} para que aparezcan en el top de la página.
    \item Comillas en español: primero angulares \guillemotleft así\guillemotright, luego inglesas, luego simples.
    \item Autores en \texttt{biblio.bib} como \emph{Apellido, Nombre} para evitar errores con nombres compuestos.
    \item Dobles llaves en títulos del \texttt{.bib} para preservar mayúsculas/minúsculas exactas.
    \item Los caracteres acentuados se escriben directamente (UTF-8), sin escapear.
    \item DOI se deja en su forma simple como campo propio; no se duplica en URL.
    \item El campo \texttt{note} del \texttt{.bib} se actualiza si la referencia es online.
    \item Cuando una sección tiene una sola subsección, agregar asterisco al comando para que no aparezca en el índice.
    \item Los diagramas apaisados —el entidad-relación y el de despliegue— van dentro de \texttt{\textbackslash begin\{landscape\}} con \texttt{width=0.95\textbackslash linewidth}. A ancho de texto vertical su tipografía interna resulta ilegible en papel. Como \texttt{landscape} fuerza un salto de página, el bloque se ubica al final de su sección y no junto al párrafo que lo menciona, para no dejar media página en blanco.
    \item Las normas argentinas se citan como \texttt{@online} en el \texttt{.bib}, con el organismo emisor entre dobles llaves como autor. Se descartó el tipo \texttt{@legislation}, que no pertenece a biblatex estándar y que el estilo \texttt{iso-authoryear} rechaza.
\end{itemize}

\subsection*{Macros disponibles}
\begin{itemize}
    \item \texttt{\textbackslash Fidel\{texto\}} — nota inline del tutor Giro Uribazo (verde).
    \item \texttt{\textbackslash Martin\{texto\}} — nota inline propia (azul).
    \item \texttt{\textbackslash pdfcomment\{texto\}} — comentario visible en el PDF.
\end{itemize}

\subsection*{Formato requerido por UADE}
\begin{itemize}
    \item Formato A4, fuente 12pt, interlineado 1.5.
    \item Referencias en rojo durante el desarrollo; cambiar a negro (\texttt{hidelinks}) en la entrega final.
    \item La carátula final se genera en la biblioteca UADE (no usar la del template).
    \item No redactar Resumen ni Abstract hasta la entrega final.
    \item Eliminar el cronograma del anexo en la entrega final.
    \item Todo contenido no propio debe estar citado, incluyendo imágenes y tablas.
\end{itemize}

\subsection*{Convenciones de redacción}
\begin{itemize}
    \item \textbf{No se usa la raya de inciso.} Decisión del autor tomada el 18/08/2026, tras comparar el registro del documento con el de la tesis de referencia Feresini e Imbriago (2026), que no emplea ese signo en ninguna de sus 141 páginas. Los incisos se resuelven con comas cuando no contienen comas internas, y con paréntesis cuando sí las contienen. La única excepción admitida es la raya como marca de dato inaplicable dentro de una celda de tabla.
    \item \textbf{Los requerimientos se enuncian en una o dos líneas.} No llevan párrafos de justificación a continuación de la tabla: el fundamento de cada decisión vive en la sección que corresponda, no en el listado.
\end{itemize}

\subsection*{Dudas pendientes}
\begin{itemize}
    \item Agregar acá las dudas que queden abiertas durante el desarrollo.
\end{itemize}
